\section{Step 74 Emit the highest selected bit first}\label{h263-writer-s04}
For the initial scalar loop, begin with source shift $n-1$. Extract one unsigned
bit, derive destination byte $p/8$ and bit position $7-(p\bmod8)$, and place it
there. Advance the local bit cursor and move to the next lower source bit until
all requested bits are consumed.

The initial empty-width return must precede the loop, so there is no source
shift of minus one. The first emitted bit belongs to the declared field, which
can contain leading zeros even when the numerical value is small.

\paragraph{Observe} Value five with width five emits \texttt{00101}, not just
\texttt{101} and not the eight-byte native representation of five.
\Needspace{8\baselineskip}
\paragraph{Check} Which source bit is emitted first for width $n>0$?
\begin{itemize}
  \item[A.] Bit $n-1$ of the right-aligned value.
  \item[B.] Always bit zero.
  \item[C.] The highest bit of the native address.
\end{itemize}
